\subsection{Fractional Exponents} 
Write a fractional exponent as $\frac{p}{q}$ in lowest terms, where $p$ is an \makeblank[1in] and $q$ is a \makeblank[1in].
\[
f(x)=a(x-h)^{p/q}+k=a\sqrt[q]{\left(x-h\right)^p}+k
\]
\begin{description}
\item[Type 1]$q$ is even, $p$ is odd, $p$ is positive; 
             $2(x-1)^{3/2}+3=2\sqrt{\left(x-1\right)^3}+3$\\
	Since we're taking an \makeblank[1in] root, \makeblank[1in] must be positive (or 0), 
	and \makeblank[1in] will be positive.
	\begin{description}
	\item[Domain] \makeblank[1in] is positive or 0, so the domain is \makeblank.
	\item[Range] \makeblank[1in] is always positive, so the range is \makeblank if \makeblanksmall is positive, or \makeblank if \makeblanksmall is negative.
	\end{description}
\item[Type 2]$q$ is even, $p$ is odd, $p$ is negative; 
             $2(x-1)^{-3/2}+3=\dfrac{2}{\sqrt{\left(x-1\right)^3}}+3$
	\begin{description}
	\item[Domain]
	\item[Range]
	\end{description}
\item[Type 3]$q$ is odd, $p$ is even, $p$ is positive; 
             $2(x-1)^{2/3}+3=2\sqrt[3]{\left(x-1\right)^2}+3$
	\begin{description}
	\item[Domain]
	\item[Range]
	\end{description}
\item[Type 4]$q$ is odd, $p$ is even, $p$ is negative; 
             $2(x-1)^{-2/3}+3=\dfrac{2}{\sqrt[3]{\left(x-1\right)^3}}+3$
	\begin{description}
	\item[Domain]
	\item[Range]
	\end{description}
\item[Type 5]$q$ is odd, $p$ is odd, $p$ is positive; 
             $2(x-1)^{3/5}+3=2\sqrt[5]{\left(x-1\right)^3}+3$
	\begin{description}
	\item[Domain]
	\item[Range]
	\end{description}
\item[Type 6]$q$ is odd, $p$ is odd, $p$ is negative; 
             $2(x-1)^{-3/5}+3=\dfrac{2}{\sqrt[5]{\left(x-1\right)^3}}+3$
	\begin{description}
	\item[Domain]
	\item[Range]
	\end{description}
\end{description}